%!TEX root = ../main.tex

\section{Experimental Setup}
\label{sec:experiments}
We measure retrieval performance on STaRK, a benchmark for entity-level retrieval over heterogeneous, text-rich KGs \cite{wu_stark_2024}.

\subsection{Benchmark}

STaRK comprises three large, heterogeneous knowledge graphs. \textbf{AMAZON} is an e-commerce graph with roughly 1M entities and 9.4M relations, constructed from Amazon metadata, reviews \cite{he_ups_2016}, and question--answer pairs \cite{mcauley_inferring_2015}. \textbf{MAG} is a scholarly graph with 1.9M entities and 39.8M relations derived from the Microsoft Academic Graph \cite{wang_microsoft_2020}. \textbf{PRIME} is a biomedical graph built from PrimeKG \cite{chandak_building_2023}, containing 129K entities and 8.1M relations. Each node is associated with text-rich attributes, making STaRK a natural testbed for hybrid retrieval over structured and textual signals. Given a query, the retriever must return a ranked list of nodes that support the answer. We report agent configuration and hyperparameters in Appendix~\ref{app:impl-details}.

\subsection{Baselines and metrics}

We compare with representative retrieval-based and agent-based baselines, focusing on methods that report results on all three graphs, as our goal is to evaluate performance consistently across different regimes and assess generality.

\slimparagraph{Retrieval-based.} \textbf{BM25} \cite{robertson_probabilistic_2009} is the same sparse, lexical retriever used for global search. We also include dense embedding retrievers that rank nodes by cosine similarity, using \textbf{ada-002} and \textbf{GritLM-7B}, an instruction-tuned 7B encoder \cite{muennighoff_generative_2025}. \textbf{mFAR}~\cite{li_multi-field_2025} is a multi-field adaptive retriever that combines keyword matching with embedding similarity to learn query-dependent weights over different node fields. \textbf{KAR}~\cite{xia_knowledge-aware_2025} augments queries with knowledge-aware expansions and applies relation-type constraints during retrieval. \textbf{MoR}~\cite{lei_mixture_2025} is a trained retriever that combines multiple retrieval objectives.

\slimparagraph{Agent-based.} \textbf{Think-on-Graph} \cite{sun_think--graph_2024} is a training-free LLM agent that iteratively expands paths in the graph using beam search. \textbf{GraphFlow} \cite{yu_can_2025} learns a policy for generating multi-hop retrieval trajectories using GFlowNets \cite{bengio_flow_2021}. \textbf{AvaTaR} \cite{wu_avatar_2024} is a tool-using agent that optimizes prompting from positive and negative trajectories.

Results for KAR, mFAR, MoR, AvaTaR, and GraphFlow are reported as in their respective papers, which evaluate on the official STaRK splits and metrics. For Think-on-Graph, we report the numbers provided in the GraphFlow study, which includes a direct comparison to Think-on-Graph under the same STaRK setup \cite{yu_can_2025,sun_think--graph_2024}.

\slimparagraph{Metrics.} We follow the STaRK protocol and report Hit@1, Hit@5, Recall@20 (R@20), and Mean Reciprocal Rank (MRR), which capture top-rank precision, coverage of the ground-truth set, and overall ranking quality. Note that Hit@5 is reported in \cref{tab:full-stark-results} in the Appendix.

\slimparagraph{Backbones.} Our primary configuration uses GPT-4.1 as the backbone LLM throughout the paper. Appendix \Cref{tab:full-stark-results} additionally reports results with GPT-4o for comparability with prior KG retrievers such as KAR and Think-on-Graph, and with GPT-5.2 to verify that the retrieval interface transfers to a newer backbone on the same splits.

\subsection{Distillation Setup}
\label{sec:distillation-setup}

For each graph, we collect teacher trajectories on the training split to distill \retrieverShort into a smaller, lower-cost model (\Cref{subsec:distill}), offering a viable alternative when under tighter compute budgets. Using GPT-4.1 as the teacher, we run \retrieverShort three times per training query with stochastic decoding (temperature $= 0.7$), producing three trajectories per query. We cap the distillation budget by subsampling up to 6{,}000 training queries per graph, yielding at most 18{,}000 trajectories per graph (full statistics in \cref{tab:traj_counts}) and summing to 94.4 million tokens. Each trajectory is limited to $T_{\max}{=}20$ steps and ends when the agent issues \textsc{finish} or reaches the step limit. We apply no trajectory filtering or rejection sampling, preserving a label-free setting. We then distill a Qwen3-8B \cite{yang_qwen3_2025} student via supervised fine-tuning with LoRA adapters \cite{hu_lora_2021}, using next-token prediction over assistant-authored tokens only. We train for one epoch with a 16{,}384-token context length using AdamW \cite{loshchilov_decoupled_2019} at learning rate $1\times 10^{-5}$, selecting checkpoints via early stopping on the official validation split. Training runs on a single NVIDIA H100 GPU and completes in approximately five hours.
